% !TeX program = xelatex
\documentclass[12pt, a4paper, titlepage, xelatex, ja=standard, jafont=hiragino-pron, twoside]{bxjsarticle}
\usepackage{xeCJK}
\usepackage{amsmath, amssymb}
\usepackage{geometry}
\usepackage{longtable}
\usepackage{booktabs}
\usepackage{caption}
\usepackage{silence}
\WarningFilter{caption}{Unknown document class}
\WarningFilter{xeCJK}{Redefining CJKfamily}
\WarningFilter{relsize}{Failed to get list}

% -- Fonts --
\setmainfont{Times New Roman}



\geometry{a4paper, margin=2.5cm}

\title{高精度界面追跡法における寄生流れの抑制と静止液滴シミュレーションの確立}
\author{}
\date{}

\begin{document}

\maketitle

\section*{Proposed Implementation and Validation Framework (CCD-CLS Approach)}

\section*{1. 背景と目的}
静止液滴計算における寄生流れは、表面張力項 $\mathbf{f}_{st} = \sigma \kappa \nabla \phi$ と圧力勾配項 $\nabla p$ の離散化誤差の不一致に起因する。本計画では、結合コンパクト差分法（CCD法）とConservative Level Set（CLS）法を用い、Balanced-Force離散化を厳密に適用することで、寄生流れを理論的限界まで抑制する計算スキームの確立を目的とする。

\section*{2. 数値計算仕様（基盤系）}
本シミュレーションで採用する「良さげな系」の具体的仕様は以下の通りとする。

\begin{center}
\begin{tabular}{|l|l|l|l|}
\hline
\textbf{要素技術} & \textbf{仕様} & \textbf{手法} & \textbf{備考} \\
\hline
空間離散化 & 6th-order CCD (Combined Compact Difference) & 速度、圧力、界面関数すべてに適用 &  \\
界面追跡 & Conservative Level Set (CLS) 法 &  & 質量保存性の確保 \\
表面張力モデル & Continuum Surface Force (CSF) + Balanced-Force & 演算子の一致を最優先 &  \\
曲率計算 & CCDによる直接2階微分評価 & $\kappa = \nabla \cdot (\frac{\nabla \phi}{|\nabla \phi|})$ &  \\
時間進行 & Fractional Step法 (Projection法) & 3rd-order Runge-Kutta 等 &  \\
格子配置 & Collocated Grid + Rhie-Chow Coupling &  & CCD-PPE（圧力ポアソン）との親和性 \\
\hline
\end{tabular}
\end{center}

\section*{3. 段階的推進ステップ（実験・調整手順）}

\subsection*{【フェーズI】幾何学的検証：曲率評価の独立テスト}
運動量方程式を解く前に、静的なレベルセット場 $\phi(\mathbf{x})$ を与え、曲率算出精度を隔離して評価する。

\textbf{実験内容}: 半径 $R$ の真円を計算領域中央に配置し、CCD法で曲率 $\kappa_{num}$ を算出。

\textbf{評価指標}: $L_2$ 誤差ノルム $E_\kappa = \sqrt{\sum (\kappa_{num} - 1/R)^2 / N}$。

\textbf{調整目標}: 格子解像度 $h$ に対して $O(h^6)$ の収束性を確認。ここで次数が出ない場合、境界条件の微分の取り扱い（Ghost Cell法等）を修正する。

\subsection*{【フェーズII】代数的検証：無粘性 Balanced-Force テスト}
粘性 $\mu$ をゼロ（または極小値）に設定し、表面張力と圧力勾配の均衡能力を試験する。

\textbf{原理}: 理論上、$\nabla p = \mathbf{f}_{st}$ が成立すれば流速 $u$ は発生しない。

\textbf{調整方法}: 圧力勾配の計算に用いるCCD演算子 $D_{ccd}$ と、表面張力項（CSF）の勾配計算に用いる $D_{ccd}$ が同一コードパスを通っているか厳密に確認。

初期流速 $u=0$ から数ステップ計算し、生じる流速 $u_{spurious}$ がマシンイプシロンに近いレベルか、あるいは $O(h^6)$ で減少するかを確認。

\subsection*{【フェーズIII】動的安定性検証：ラプラス数感度分析}
粘性を導入し、物理的な散逸がある状態での寄生流れの定常値を評価する。

\textbf{実験パラメータ}: $La = \sigma \rho D / \mu^2$ を $10 \sim 10^5$ の範囲で振る。

\textbf{評価指標}: キャピラリ数 $Ca = \mu U_{max} / \sigma$ の時間推移。

\textbf{調整方法}: $Ca$ が発散する場合、界面厚さパラメータ $\epsilon$ を $0.5h \sim 1.5h$ の範囲で調整し、数値振動（オーバーシュート）を抑制する。

\section*{4. 寄生流れを抑制するための調整ガイドライン}

\subsection*{1. 界面スムージングの最適化}
鋭すぎる界面はCCD法の高次精度を台無しにする。以下のヘビサイド関数 $H(\phi)$ を推奨する。
$$H(\phi) = \begin{cases}
0 & (\phi < -\epsilon) \\
\frac{1}{2} \left[1 + \frac{\phi}{\epsilon} + \frac{1}{\pi} \sin\left(\frac{\pi \phi}{\epsilon}\right)\right] & (|\phi| \le \epsilon) \\
1 & (\phi > \epsilon)
\end{cases}$$
※この導関数（デルタ関数）が $C^1$ 連続であることを利用し、表面張力分布を滑らかにする。

\subsection*{2. PPE（圧力ポアソン方程式）の収束条件}
寄生流れは「打ち消し残差」である。仕様: 収束判定値を $10^{-10}$ 程度に厳格化。手法: BiCGSTAB法などの反復解法を用いつつ、CCDの三対角構造を活かした前処理を検討。

\subsection*{3. 時間刻み幅 $\Delta t$ の制約}
表面張力に起因する毛管波の安定性条件（Brackbillらの条件）を厳守する。
$$\Delta t < \sqrt{\frac{\rho h^3}{2\pi\sigma}}$$
CCD法は高次精度だが、時間方向の安定限界は物理現象に縛られるため、CFL条件よりもこの毛管波条件が厳しくなることが多い。

\section*{5. 結論と期待される成果}
以上の手順により、寄生流れの発生源を「幾何誤差（曲率）」「代数誤差（演算子の不一致）」「物理不安定（界面厚さと時間刻み）」に切り分けて対策することが可能となる。最終的に、高 $La$ 数条件下でも液滴形状が維持され、$Ca < 10^{-5}$ 程度の低ノイズな静止液滴シミュレーションの実現が期待される。

\end{document}